\documentclass[11pt]{article}

\usepackage[margin=1in]{geometry}
\usepackage{amsmath}
\usepackage{amssymb}
\usepackage{booktabs}
\usepackage{array}
\usepackage[hidelinks]{hyperref}

\newcommand{\softmax}{\operatorname{softmax}}
\newcommand{\Concat}{\operatorname{Concat}}

\title{\textbf{Exercise 03: Multi-Head Self-Attention}}
\author{}
\date{}

\begin{document}
\maketitle

\begin{center}
\emph{A numerical two-head self-attention example with independent projections.}
\end{center}

\section*{Setup}

Use the row-token input
\[
X=
\begin{bmatrix}
1&0&1&0\\
0&1&1&1\\
1&1&0&1
\end{bmatrix}
\in\mathbb{R}^{3\times4}.
\]
There are $H=2$ heads with $d_k=d_v=2$ per head. The projection matrices are fixed toy values; nothing is trained in this exercise.

\section*{Head 1}

Head 1 uses
\[
W_Q^{(1)}=W_K^{(1)}=W_V^{(1)}=
\begin{bmatrix}
1&0\\
0&1\\
0&0\\
0&0
\end{bmatrix}.
\]
Therefore
\[
Q_1=K_1=V_1=
\begin{bmatrix}
1&0\\
0&1\\
1&1
\end{bmatrix}.
\]

The scaled score matrix is
\[
\frac{Q_1K_1^\top}{\sqrt{2}}
\approx
\begin{bmatrix}
0.7071&0&0.7071\\
0&0.7071&0.7071\\
0.7071&0.7071&1.4142
\end{bmatrix}.
\]
Applying row-wise softmax gives
\[
A_1\approx
\begin{bmatrix}
0.4011&0.1978&0.4011\\
0.1978&0.4011&0.4011\\
0.2483&0.2483&0.5035
\end{bmatrix}.
\]
The first head output is
\[
O_1=A_1V_1
\approx
\begin{bmatrix}
0.8022&0.5989\\
0.5989&0.8022\\
0.7517&0.7517
\end{bmatrix}.
\]

\section*{Head 2}

Head 2 receives the same $X$ but uses different projections:
\[
W_Q^{(2)}=
\begin{bmatrix}
0&0\\
0&0\\
1&0\\
0&1
\end{bmatrix},\quad
W_K^{(2)}=
\begin{bmatrix}
0&0\\
0&0\\
0&1\\
1&0
\end{bmatrix},\quad
W_V^{(2)}=
\begin{bmatrix}
0&0\\
0&0\\
1&0\\
0&1
\end{bmatrix}.
\]
This gives
\[
Q_2=
\begin{bmatrix}
1&0\\
1&1\\
0&1
\end{bmatrix},\qquad
K_2=
\begin{bmatrix}
0&1\\
1&1\\
1&0
\end{bmatrix},\qquad
V_2=
\begin{bmatrix}
1&0\\
1&1\\
0&1
\end{bmatrix}.
\]

The scaled score matrix is
\[
\frac{Q_2K_2^\top}{\sqrt{2}}
\approx
\begin{bmatrix}
0&0.7071&0.7071\\
0.7071&1.4142&0.7071\\
0.7071&0.7071&0
\end{bmatrix}.
\]
Hence
\[
A_2\approx
\begin{bmatrix}
0.1978&0.4011&0.4011\\
0.2483&0.5035&0.2483\\
0.4011&0.4011&0.1978
\end{bmatrix},
\]
and
\[
O_2=A_2V_2
\approx
\begin{bmatrix}
0.5989&0.8022\\
0.7517&0.7517\\
0.8022&0.5989
\end{bmatrix}.
\]

\section*{Why the heads differ}

The two heads receive the same input matrix $X$, but they do not see the same projected geometry. Their independent $W_Q$, $W_K$, and $W_V$ matrices produce different query-key dot products and therefore different attention distributions:
\[
A_1\neq A_2.
\]

This controlled example demonstrates only that independent projections can produce distinct attention patterns. The projection matrices are hand-designed rather than learned, so the exercise does \emph{not} assign semantic roles such as ``relational'' or ``anchoring'' to either head.

\section*{Concatenation}

Concatenate the head outputs along the feature dimension:
\[
O_{\mathrm{concat}}=\Concat(O_1,O_2)
\approx
\begin{bmatrix}
0.8022&0.5989&0.5989&0.8022\\
0.5989&0.8022&0.7517&0.7517\\
0.7517&0.7517&0.8022&0.5989
\end{bmatrix}.
\]
Thus
\[
O_{\mathrm{concat}}\in\mathbb{R}^{3\times4}.
\]

\section*{Output projection}

For this toy exercise,
\[
W_O=I_4.
\]
Therefore
\[
O_{\mathrm{MHA}}=O_{\mathrm{concat}}W_O=O_{\mathrm{concat}}.
\]
The identity choice is only a simplification that keeps the concatenated features visible. In a trained Transformer, $W_O$ is learned and mixes information across head features.

\section*{Causal extension}

If the same two-head layer were causal, the same causal mask $M$ would be added inside each head before softmax:
\[
A_{h,\mathrm{causal}}
=
\softmax\left(
\frac{Q_hK_h^\top}{\sqrt{d_k}}+M
\right).
\]
The visibility constraint is shared, while each head still has its own projected logits.

\section*{Pipeline summary}

For each head $h$,
\[
Q_h=XW_Q^{(h)},\qquad
K_h=XW_K^{(h)},\qquad
V_h=XW_V^{(h)},
\]
\[
\mathrm{head}_h
=
\softmax\left(\frac{Q_hK_h^\top}{\sqrt{d_k}}\right)V_h.
\]
Then
\[
\boxed{
O_{\mathrm{MHA}}
=
\Concat(\mathrm{head}_1,\ldots,\mathrm{head}_H)W_O
}.
\]

\section*{Numerical verification}

The companion \texttt{answer\_key.py} reproduces the two attention matrices, both head outputs, the concatenation, and the final identity output projection with deterministic assertions.

\end{document}
